% GENERATED FILE. Edit canonical lessons, then run npm run generate:books.
% canonical-source-hash: fnv1a64:45fa9ba06041f1f7
% canonical-lessons: ES-C679-comisaria, ES-C679-embajada, ES-C679-estanco, ES-C679-sello, ES-C679-aseo, ES-R679-repaso-672-679

\chapter{Servicios de la ciudad}
\label{ch:es-servicios-de-la-ciudad}
\hlchaptermodality{679}

\begin{chapteropening}
\textbf{By the end of this chapter:} \emph{I can find the police station, the embassy, the tobacconist, stamps and the toilets.}

Complete the last lesson of chapter 679: I can find the police station, the embassy, the tobacconist, stamps and the toilets.
\end{chapteropening}

\section[la comisaría]{la comisaría — the police station}

\label{lesson:ES-C679-comisaria}



\begin{quote}
Before the new one: say the Spanish for the tower, then the Spanish for the gallery.
\end{quote}

\subsection*{You'll want to know: la comisaría}
\textbf{la comisaría} — \textquotedblleft{}the police station\textquotedblright{}.

\begin{quote}
La comisaría está al lado del ayuntamiento.
\end{quote}

\begin{quote}
\emph{The police station is next to the town hall.}
\end{quote}

In Latin America also \emph{la estación de policía}.

\begin{cousinweb}
\emph{comisaría} is built on \emph{comisario} \textquotedblleft{}commissioner\textquotedblright{}, from Medieval Latin \emph{commissārius} \textquotedblleft{}one entrusted\textquotedblright{}, from \emph{committere} \textquotedblleft{}to entrust\textquotedblright{}.
\end{cousinweb}

\begin{grammarlens}[title={What you've built}]
One more word for telling, explaining and describing at B1.
\end{grammarlens}

\subsection*{Your turn}
\begin{itemize}
\raggedright
  \item \emph{Say it:} \emph{la comisaría}
  \item \emph{Say it:} \emph{la comisaría}, once more
  \item \emph{Say it:} a sentence of your own with \emph{la comisaría}
  \item \emph{Recall:} say the Spanish for the tower, then the Spanish for the gallery, then say \emph{la comisaría} again
\end{itemize}

\begin{tcolorbox}[breakable,colback=teal!4,colframe=teal!35!black,title={Before you move on}]
What does la comisaría mean? (\textquotedblleft{}The police station\textquotedblright{}.) Say it once more.
\end{tcolorbox}
\section[la embajada]{la embajada — the embassy}

\label{lesson:ES-C679-embajada}



\begin{quote}
Before the new one: say the Spanish for the gallery, then the Spanish for the police station.
\end{quote}

\subsection*{You'll want to know: la embajada}
\textbf{la embajada} — \textquotedblleft{}the embassy\textquotedblright{}.

\begin{quote}
Perdí el pasaporte y llamé a la embajada.
\end{quote}

\begin{quote}
\emph{I lost my passport and phoned the embassy.}
\end{quote}

The person is \emph{el embajador}, \emph{la embajadora}.

\begin{cousinweb}
\emph{embajada} came from Occitan \emph{ambaissada} \textquotedblleft{}embassy\textquotedblright{} (EtymDB), from Late Latin \emph{ambactia}, built on Gaulish \emph{ambactos} \textquotedblleft{}a servant, a messenger\textquotedblright{}.
\end{cousinweb}

\begin{grammarlens}[title={What you've built}]
One more word for telling, explaining and describing at B1.
\end{grammarlens}

\subsection*{Your turn}
\begin{itemize}
\raggedright
  \item \emph{Say it:} \emph{la embajada}
  \item \emph{Say it:} \emph{la embajada}, once more
  \item \emph{Say it:} a sentence of your own with \emph{la embajada}
  \item \emph{Recall:} say the Spanish for the gallery, then the Spanish for the police station, then say \emph{la embajada} again
\end{itemize}

\begin{tcolorbox}[breakable,colback=teal!4,colframe=teal!35!black,title={Before you move on}]
What does la embajada mean? (\textquotedblleft{}The embassy\textquotedblright{}.) Say it once more.
\end{tcolorbox}
\section[el estanco]{el estanco — the tobacconist’s (Spain)}

\label{lesson:ES-C679-estanco}



\begin{quote}
Before the new one: say the Spanish for the police station, then the Spanish for the embassy.
\end{quote}

\subsection*{You'll want to know: el estanco}
\textbf{el estanco} — \textquotedblleft{}the tobacconist’s (Spain)\textquotedblright{}.

\begin{quote}
El estanco de la plaza abre los domingos.
\end{quote}

\begin{quote}
\emph{The tobacconist on the square opens on Sundays.}
\end{quote}

In Spain \emph{el estanco} sells tobacco, and stamps too.

\begin{cousinweb}
\emph{estanco} is built on \emph{estancar} \textquotedblleft{}to hold back\textquotedblright{}, which EtymDB traces to Latin \emph{stagnum} \textquotedblleft{}standing water\textquotedblright{}; the shop once sold goods the state kept for itself.
\end{cousinweb}

\begin{grammarlens}[title={What you've built}]
One more word for telling, explaining and describing at B1.
\end{grammarlens}

\subsection*{Your turn}
\begin{itemize}
\raggedright
  \item \emph{Say it:} \emph{el estanco}
  \item \emph{Say it:} \emph{el estanco}, once more
  \item \emph{Say it:} a sentence of your own with \emph{el estanco}
  \item \emph{Recall:} say the Spanish for the police station, then the Spanish for the embassy, then say \emph{el estanco} again
\end{itemize}

\begin{tcolorbox}[breakable,colback=teal!4,colframe=teal!35!black,title={Before you move on}]
What does el estanco mean? (\textquotedblleft{}The tobacconist’s (Spain)\textquotedblright{}.) Say it once more.
\end{tcolorbox}
\section[el sello]{el sello — the stamp; the seal}

\label{lesson:ES-C679-sello}



\begin{quote}
Before the new one: say the Spanish for the embassy, then the Spanish for the tobacconist’s.
\end{quote}

\subsection*{You'll want to know: el sello}
\textbf{el sello} — \textquotedblleft{}the stamp; the seal\textquotedblright{}.

\begin{quote}
Necesito un sello para esta carta.
\end{quote}

\begin{quote}
\emph{I need a stamp for this letter.}
\end{quote}

In Latin America \emph{la estampilla}, and in Mexico \emph{el timbre}.

\begin{cousinweb}
\emph{sello} is inherited from Latin \emph{sigillum} \textquotedblleft{}a little image, a seal\textquotedblright{}, a little \emph{signum} \textquotedblleft{}sign\textquotedblright{}.
\end{cousinweb}

\begin{grammarlens}[title={What you've built}]
One more word for telling, explaining and describing at B1.
\end{grammarlens}

\subsection*{Your turn}
\begin{itemize}
\raggedright
  \item \emph{Say it:} \emph{el sello}
  \item \emph{Say it:} \emph{el sello}, once more
  \item \emph{Say it:} a sentence of your own with \emph{el sello}
  \item \emph{Recall:} say the Spanish for the embassy, then the Spanish for the tobacconist’s, then say \emph{el sello} again
\end{itemize}

\begin{tcolorbox}[breakable,colback=teal!4,colframe=teal!35!black,title={Before you move on}]
What does el sello mean? (\textquotedblleft{}The stamp; the seal\textquotedblright{}.) Say it once more.
\end{tcolorbox}
\section[el aseo]{el aseo — the toilet, the restroom (Spain); cleanliness}

\label{lesson:ES-C679-aseo}



\begin{quote}
Before the new one: say the Spanish for the tobacconist’s, then the Spanish for the stamp.
\end{quote}

\subsection*{You'll want to know: el aseo}
\textbf{el aseo} — \textquotedblleft{}the toilet, the restroom (Spain); cleanliness\textquotedblright{}.

\begin{quote}
¿Dónde están los aseos?
\end{quote}

\begin{quote}
\emph{Where are the toilets?}
\end{quote}

\emph{Los aseos} on a sign in Spain; Latin America says \emph{el baño}.

\begin{cousinweb}
\emph{aseo} is built on \emph{asear} \textquotedblleft{}to tidy, to clean\textquotedblright{} (EtymDB). The origin of \emph{asear} is not settled.
\end{cousinweb}

\begin{grammarlens}[title={What you've built}]
One more word for telling, explaining and describing at B1.
\end{grammarlens}

\subsection*{Your turn}
\begin{itemize}
\raggedright
  \item \emph{Say it:} \emph{el aseo}
  \item \emph{Say it:} \emph{el aseo}, once more
  \item \emph{Say it:} a sentence of your own with \emph{el aseo}
  \item \emph{Recall:} say the Spanish for the tobacconist’s, then the Spanish for the stamp, then say \emph{el aseo} again
\end{itemize}

\begin{tcolorbox}[breakable,colback=teal!4,colframe=teal!35!black,title={Before you move on}]
What does el aseo mean? (\textquotedblleft{}The toilet, the restroom (Spain); cleanliness\textquotedblright{}.) Say it once more.
\end{tcolorbox}
\section[(dialogue)]{Repaso de los capítulos 672–679}

\label{lesson:ES-R679-repaso-672-679}



\begin{quote}
Say the Spanish for the signpost and the sign.
\end{quote}

\subsection*{Your turn}
Say these aloud:

\begin{itemize}
\raggedright
  \item the signpost, the sign, the arrow, the signs, the diversion
  \item to head, to shorten, to go back, to go round, to go along the edge of
  \item to go for a walk, the walk, to wander the streets, to travel around, the route
  \item the outskirts, the block, the alley, the arcade, the paving stone
  \item tourism, the viewpoint, tourist, the leaflet, to haggle
  \item the monument, the cathedral, the castle, the palace, the city wall
  \item to visit, the country chapel, the statue, the tower, the gallery
  \item the police station, the embassy, the tobacconist’s, the stamp, the toilet
\end{itemize}

\begin{tcolorbox}[breakable,colback=teal!4,colframe=teal!35!black,title={Before you move on}]
The signpost? (\textbf{el indicador}.) The sign? (\textbf{el letrero}.) The arrow? (\textbf{la flecha}.) The signs? (\textbf{la señalización}.) The diversion? (\textbf{el desvío}.) To head? (\textbf{dirigirse}.) To shorten? (\textbf{acortar}.) To go back? (\textbf{retroceder}.) To go round? (\textbf{rodear}.) To go along the edge of? (\textbf{bordear}.) To go for a walk? (\textbf{pasear}.) The walk? (\textbf{el paseo}.) To wander the streets? (\textbf{callejear}.) To travel around? (\textbf{recorrer}.) The route? (\textbf{el recorrido}.) The outskirts? (\textbf{las afueras}.) The block? (\textbf{la cuadra}.) The alley? (\textbf{el callejón}.) The arcade? (\textbf{el soportal}.) The paving stone? (\textbf{el adoquín}.) Tourism? (\textbf{el turismo}.) The viewpoint? (\textbf{el mirador}.) Tourist? (\textbf{turístico}.) The leaflet? (\textbf{el folleto}.) To haggle? (\textbf{regatear}.) The monument? (\textbf{el monumento}.) The cathedral? (\textbf{la catedral}.) The castle? (\textbf{el castillo}.) The palace? (\textbf{el palacio}.) The city wall? (\textbf{la muralla}.) To visit? (\textbf{visitar}.) The country chapel? (\textbf{la ermita}.) The statue? (\textbf{la estatua}.) The tower? (\textbf{la torre}.) The gallery? (\textbf{la galería}.) The police station? (\textbf{la comisaría}.) The embassy? (\textbf{la embajada}.) The tobacconist’s? (\textbf{el estanco}.) The stamp? (\textbf{el sello}.) The toilet? (\textbf{el aseo}.)
\end{tcolorbox}
